% GENERATED FILE. Edit canonical lessons, then run npm run generate:books.
% canonical-source-hash: fnv1a64:8f7b3fc017e87d44
% canonical-lessons: ZH-C171-dakai, ZH-C171-kongtiao, ZH-C171-yaokongqi, ZH-C171-an, ZH-C171-peng

\chapter{Switching Things On}
\label{ch:zh-switching-things-on}
\hlchaptermodality{171}

\begin{chapteropening}
\textbf{By the end of this chapter:} \emph{I can open things, put the air conditioning on, use the remote, press buttons and say what not to touch.}

Complete the last lesson of chapter 171: I can open things, put the air conditioning on, use the remote, press buttons and say what not to touch.
\end{chapteropening}

\section[\texorpdfstring{\zh{打开} (dǎkāi)}{打开 (dǎkāi)}]{\zh{打开} (dǎkāi) — to open; to switch on}

\label{lesson:ZH-C171-dakai}



\begin{quote}
Before the new one: say the Mandarin for a fork, then the Mandarin for a microwave oven.
\end{quote}

\subsection*{You'll want to know: \zh{打开}}
\textbf{\zh{打开}} — \emph{dǎkāi} — \textquotedblleft{}to open; to switch on\textquotedblright{}.

\begin{quote}
\textbf{\zh{你去打开书房的门}} \emph{nǐ qù dǎkāi shūfáng de mén} — Go and open the study door.
\end{quote}

\begin{cousinweb}
\textbf{\zh{打开}} (\emph{dǎkāi}) is built from \zh{打} 'to strike' and \zh{开} 'to open', both inherited from Middle Chinese, \emph{taengX} and \emph{khoj} (Guangyun, in Baxter's transcription), though dǎ is not the sound \emph{taengX} would regularly give.
\end{cousinweb}

\begin{grammarlens}[title={What you've built}]
One more word for saying what I do.
\end{grammarlens}

\subsection*{Your turn}
\begin{itemize}
\raggedright
  \item \emph{Say it:} \emph{dǎkāi}
  \item \emph{Say it:} \emph{dǎkāi}, once more
  \item \emph{Say it:} \emph{nǐ qù dǎkāi shūfáng de mén} — Go and open the study door.
  \item \emph{Recall:} say the Mandarin for a fork, then the Mandarin for a microwave oven, then say \emph{dǎkāi} again
\end{itemize}

\begin{tcolorbox}[breakable,colback=teal!4,colframe=teal!35!black,title={Before you move on}]
What does \zh{打开} mean? (\textquotedblleft{}To open; to switch on\textquotedblright{}.) Say it once more.
\end{tcolorbox}
\section[kōngtiáo]{kōngtiáo — air conditioner; air conditioning}

\label{lesson:ZH-C171-kongtiao}



\begin{quote}
Before the new one: say the Mandarin for a microwave oven, then the Mandarin for to switch on.
\end{quote}

\subsection*{You'll want to know: kōngtiáo}
\emph{kōngtiáo} — \textquotedblleft{}air conditioner; air conditioning\textquotedblright{}.

\begin{quote}
\emph{jīntiān tài rè le wǒmen kāi kōngtiáo} — It's far too hot today: let's put the air conditioner on.
\end{quote}

Its characters come later; for now it is a word for the ear.

\begin{cousinweb}
\emph{kōngtiáo} is built from \emph{kōng} 'empty', here the air, and \emph{tiáo} 'to adjust': a short name for the machine that adjusts the air. The first is inherited from Middle Chinese \emph{khuwng} (Guangyun, in Baxter's transcription), glossed 'empty'; the second is listed there as \emph{dew}, glossed 'to harmonize'.
\end{cousinweb}

\begin{grammarlens}[title={What you've built}]
One more word for saying what I do.
\end{grammarlens}

\subsection*{Your turn}
\begin{itemize}
\raggedright
  \item \emph{Say it:} \emph{kōngtiáo}
  \item \emph{Say it:} \emph{kōngtiáo}, once more
  \item \emph{Say it:} \emph{jīntiān tài rè le wǒmen kāi kōngtiáo} — It's far too hot today: let's put the air conditioner on.
  \item \emph{Recall:} say the Mandarin for a microwave oven, then the Mandarin for to switch on, then say \emph{kōngtiáo} again
\end{itemize}

\begin{tcolorbox}[breakable,colback=teal!4,colframe=teal!35!black,title={Before you move on}]
What does kōngtiáo mean? (\textquotedblleft{}Air conditioner; air conditioning\textquotedblright{}.) Say it once more.
\end{tcolorbox}
\section[yáokòngqì]{yáokòngqì — remote control}

\label{lesson:ZH-C171-yaokongqi}



\begin{quote}
Before the new one: say the Mandarin for to switch on, then the Mandarin for an air conditioner.
\end{quote}

\subsection*{You'll want to know: yáokòngqì}
\emph{yáokòngqì} — \textquotedblleft{}remote control\textquotedblright{}.

\begin{quote}
\emph{yáokòngqì zài nǎr} — Where's the remote?
\end{quote}

Its characters come later; for now it is a word for the ear.

\begin{cousinweb}
\emph{yáokòngqì} is built from \emph{yáo} 'far', \emph{kòng} 'to control' and \emph{qì} 'device'. The last is inherited from Middle Chinese \emph{khijH} (Guangyun, in Baxter's transcription), glossed 'a vessel'; the first two are listed there as \emph{yew}, glossed 'far', and \emph{khuwngH}, glossed 'to pull': a device that pulls from far away.
\end{cousinweb}

\begin{grammarlens}[title={What you've built}]
One more word for saying what I do.
\end{grammarlens}

\subsection*{Your turn}
\begin{itemize}
\raggedright
  \item \emph{Say it:} \emph{yáokòngqì}
  \item \emph{Say it:} \emph{yáokòngqì}, once more
  \item \emph{Say it:} \emph{yáokòngqì zài nǎr} — Where's the remote?
  \item \emph{Recall:} say the Mandarin for to switch on, then the Mandarin for an air conditioner, then say \emph{yáokòngqì} again
\end{itemize}

\begin{tcolorbox}[breakable,colback=teal!4,colframe=teal!35!black,title={Before you move on}]
What does yáokòngqì mean? (\textquotedblleft{}Remote control\textquotedblright{}.) Say it once more.
\end{tcolorbox}
\section[àn]{àn — to press; to push down}

\label{lesson:ZH-C171-an}



\begin{quote}
Before the new one: say the Mandarin for an air conditioner, then the Mandarin for a remote control.
\end{quote}

\subsection*{You'll want to know: àn}
\emph{àn} — \textquotedblleft{}to press; to push down\textquotedblright{}.

\begin{quote}
\emph{qǐng àn zhège} — Please press this one.
\end{quote}

Its characters come later; for now it is a word for the ear.

\begin{cousinweb}
\emph{àn} is inherited from Middle Chinese \emph{'anH} (Guangyun, in Baxter's transcription), glossed 'to press down' and 'to stop'.
\end{cousinweb}

\begin{grammarlens}[title={What you've built}]
One more word for saying what I do.
\end{grammarlens}

\subsection*{Your turn}
\begin{itemize}
\raggedright
  \item \emph{Say it:} \emph{àn}
  \item \emph{Say it:} \emph{àn}, once more
  \item \emph{Say it:} \emph{qǐng àn zhège} — Please press this one.
  \item \emph{Recall:} say the Mandarin for an air conditioner, then the Mandarin for a remote control, then say \emph{àn} again
\end{itemize}

\begin{tcolorbox}[breakable,colback=teal!4,colframe=teal!35!black,title={Before you move on}]
What does àn mean? (\textquotedblleft{}To press; to push down\textquotedblright{}.) Say it once more.
\end{tcolorbox}
\section[pèng]{pèng — to touch; to bump into}

\label{lesson:ZH-C171-peng}



\begin{quote}
Before the new one: say the Mandarin for a remote control, then the Mandarin for to press a button.
\end{quote}

\subsection*{You'll want to know: pèng}
\emph{pèng} — \textquotedblleft{}to touch; to bump into\textquotedblright{}.

\begin{quote}
\emph{bié pèng zhège hěn rè} — Don't touch that: it's hot.
\end{quote}

Its characters come later; for now it is a word for the ear.

\begin{cousinweb}
\emph{pèng} No dictionary we use records an older form, so its history is not settled.
\end{cousinweb}

\begin{grammarlens}[title={What you've built}]
One more word for saying what I do.
\end{grammarlens}

\subsection*{Your turn}
\begin{itemize}
\raggedright
  \item \emph{Say it:} \emph{pèng}
  \item \emph{Say it:} \emph{pèng}, once more
  \item \emph{Say it:} \emph{bié pèng zhège hěn rè} — Don't touch that: it's hot.
  \item \emph{Recall:} say the Mandarin for a remote control, then the Mandarin for to press a button, then say \emph{pèng} again
\end{itemize}

\begin{tcolorbox}[breakable,colback=teal!4,colframe=teal!35!black,title={Before you move on}]
What does pèng mean? (\textquotedblleft{}To touch; to bump into\textquotedblright{}.) Say it once more.
\end{tcolorbox}
